% =========================================================
% 2.1 SAFETIPIN
% El \section{} de esta sección está en el chapter.tex del capítulo;
% sus referencias van en seccion2_1.bib (misma carpeta).
% =========================================================

SafetiPin es una iniciativa lanzada en diciembre de 2013 en India por Kalpana Viswanath. Las dos premisas principales que sostiene la aplicación son que la recolección de información a gran escala puede guiar un cambio y que la seguridad se garantizará mientras más personas se involucren en el problema\cite{viswanath2015safetipin}.

La idea de la aplicación es usar el concepto de \enquote{safety audit}, un método para identificar lugares públicos específicos como inseguros, así como los factores que los hacen percibirse de tal manera\cite{whitzman2014partnerships}. Si bien el concepto de las auditorías de seguridad está pensado originalmente para las mujeres, la tesis central es que, al hacer los espacios más seguros para ellas, se hacen más seguros para todos\cite{viswanath2015safetipin}.

De esta manera, se plantea que las personas se involucren en el proceso de creación de espacios más seguros, calificándolos a lo largo del planeta con una rúbrica de auditoría de seguridad que va de 1 a 4, como se muestra en el cuadro \ref{cuadro:safeti_pin_rubric}. Esto, combinado con el trabajo de auditores y voluntarios, conforma una base de datos centralizada que permite a otras personas, gobiernos y asociaciones sumarse para recabar información y comprender los factores que generan inseguridad.

SafetiPin es, sin duda, el aplicativo con más trayectoria y respaldo detrás. Como parte de una iniciativa de alcance internacional con presencia en más de 70 ciudades y alrededor de 600,000 auditorías realizadas, representa un caso de éxito consolidado del crowdsourcing como medio para recabar información. Su motor central es la rúbrica de auditoría de seguridad, que constituye, en esencia, una aplicación práctica de conceptos de criminología ambiental, principalmente de la prevención del crimen mediante el diseño urbano.

\begin{table}[h]
\centering
\scriptsize
\renewcommand{\arraystretch}{1.5}
\begin{tabular}{|p{2cm}|p{2.8cm}|p{2.8cm}|p{2.8cm}|p{2.8cm}|}
\hline
& \textbf{1} & \textbf{2} & \textbf{3} & \textbf{4} \\ \hline
\textbf{Luz (noche)} & \textbf{Ninguna.} Sin luces de la calle u otras. & \textbf{Poca.} Se ven luces, pero hay baja visibilidad. & \textbf{Suficiente.} La iluminación basta. & \textbf{Brillante.} Toda el área está muy iluminada. \\ \hline
\textbf{Apertura} & \textbf{Cerrada.} Muchos puntos ciegos y sin línea de visión clara. & \textbf{Parcialmente abierta.} Se puede ver un poco hacia adelante y alrededor. & \textbf{Mayormente abierta.} Se puede ver en la mayoría de las direcciones. & \textbf{Completamente abierta.} Se puede ver claramente en todas las direcciones. \\ \hline
\textbf{Visibilidad} & \textbf{Nula.} Ninguna ventana o entrada de tiendas o residencias da a este punto. & \textbf{Escasa.} Menos de 5 ventanas o entradas dan a este punto. & \textbf{Moderada.} Menos de 10 ventanas o entradas dan a este punto. & \textbf{Alta.} Más de 10 ventanas o entradas dan a este punto. \\ \hline
\textbf{Personas} & \textbf{Desierto.} Nadie a la vista. & \textbf{Pocas personas.} Menos de 10 personas a la vista. & \textbf{Concurrencia moderada.} Más de 10 personas visibles. & \textbf{Abarrotado.} Muchas personas a una distancia en la que se pueden tocar. \\ \hline
\textbf{Seguridad} & \textbf{Ninguna.} Sin guardias ni policía visibles en los alrededores. & \textbf{Mínima.} Algo de seguridad privada visible en los alrededores pero no cerca. & \textbf{Moderada.} Seguridad privada a una distancia desde donde pueden escuchar un llamado. & \textbf{Alta.} Policía / seguridad confiable a una distancia desde donde pueden escuchar un llamado. \\ \hline
\textbf{Vía peatonal} & \textbf{Ninguna.} No hay camino para caminar disponible. & \textbf{Mala.} Existe un camino pero en muy malas condiciones. & \textbf{Regular.} Se puede caminar pero no correr. & \textbf{Buena.} Fácil de caminar rápido o correr. \\ \hline
\textbf{Transporte Público} & \textbf{No disponible.} Sin parada de metro, autobús, auto/rickshaw a menos de 10 minutos a pie. & \textbf{Distante.} Parada de metro, autobús, auto/rickshaw entre 5 y 10 minutos a pie. & \textbf{Cercano.} Parada de metro, autobús, auto/rickshaw entre 2 y 5 minutos a pie. & \textbf{Muy cerca.} Parada de metro, autobús, auto/rickshaw disponible a menos de 2 minutos a pie. \\ \hline
\textbf{Diversidad de Género} & \textbf{No diversa.} Nadie a la vista, o solo hombres. & \textbf{Algo diversa.} Mayoría de hombres, muy pocas mujeres o niños. & \textbf{Bastante diversa.} Algunas mujeres y niños. & \textbf{Diversa.} Equilibrio de todos los géneros o más mujeres y niños. \\ \hline
\textbf{Sensación} & \textbf{Aterradora.} Nunca me aventuraría aquí sin escolta suficiente. & \textbf{Incómoda.} La evitaré siempre que sea posible. & \textbf{Aceptable.} Tomaré otras rutas disponibles y mejores cuando sea posible. & \textbf{Cómoda.} Puedo tomar esta ruta incluso de noche. \\ \hline
\end{tabular}
\caption{Rúbrica de auditoría de seguridad}
\label{cuadro:safeti_pin_rubric}
{\footnotesize Tabla traducida del artículo \enquote{SafetiPin: an innovative mobile app to collect data on women's safety in Indian cities} \cite{viswanath2015safetipin}.}
\end{table}
